\section{Introduction}\label{sec:intro}

\subsection{Background}
Office buildings now carry dense sensor networks. Building-management systems log electricity use, supply- and return-air temperatures, fan speeds, indoor temperatures and on-site weather every few minutes. Facility managers use these streams to predict heating, ventilation and air-conditioning (HVAC) energy, to detect faults and to plan demand response. Buildings account for a large share of electricity use, and HVAC is usually the largest controllable end use, so accurate HVAC prediction has direct cost and carbon value \citep{amasyali2018review}.

Real sensor data are rarely complete. Sensors fail, controllers drop off the network, loggers drift, and some points are never installed because they are expensive or intrusive. The public three-year dataset of LBNL Building 59 \citep{luo2022dataset}, used in this project, shows the problem clearly. Even after the authors' curation, the supply-air temperature is missing for 7.9\% of the hours, the HVAC meter for 4.6\% and the Wi-Fi occupancy counts for more than half of the period (Section~\ref{sec:data}).

Purely data-driven models learn statistical associations between inputs and energy. When an input disappears at run time, the model sees a pattern it was never trained on, and its prediction can become both inaccurate and physically implausible. Physics-informed machine learning (PIML) adds prior physical knowledge, such as an energy balance, to the architecture or to the loss function \citep{karniadakis2021piml,raissi2019pinn}. For buildings, a first-order resistance--capacitance (RC) model links indoor temperature change to outdoor temperature, internal gains, solar gains and HVAC power \citep{bacher2011greybox}. Such a prior can constrain a neural network when parts of the data are missing.

\subsection{Problem statement and baseline}
The principal baseline of this project is \citet{mahajan2024bnn}. That study predicts hourly HVAC energy of LBNL Building 59 from weather, indoor temperature, other electricity loads and calendar variables with a Bayesian neural network (BNN). It trains on about 2.5 years and tests on the last six months of 2020, and reports a test RMSE of 9.65\,kWh, an MAE of 7.06\,kWh and a MAPE of 0.35 for the BNN (11.86, 8.21 and 0.45 for an MC-dropout LSTM). This is the closest published work: same building, same target, same test window.

The baseline has three gaps that this project addresses:
\begin{enumerate}[nosep]
\item \textbf{Complete data only.} Missing values are forward-filled before modelling and the models are never tested with missing sensors. Its robustness to sensor loss is unknown.
\item \textbf{No physical knowledge.} The BNN is a generic multilayer perceptron. Nothing ties its predictions to the heat balance of the building.
\item \textbf{No temporal context and no missingness signal.} Each hour is predicted from that hour's readings only, and an imputed value looks the same to the model as a measured one.
\end{enumerate}
Physics-informed building models exist \citep{gokhale2022pinn,dinatale2022pcnn,chen2023physcon}, and so do studies of missing building data \citep{kim2023imputation,lee2024missing,liguori2024opening}. They are, however, rarely brought together: no study found in this review tests whether a physics prior protects an operational HVAC energy predictor when its input sensors fail (Section~\ref{sec:gap}).

\subsection{Research question and objectives}
\textbf{Research question.} \emph{To what extent does a mask-aware, physics-informed recurrent network improve hourly HVAC energy prediction over conventional machine-learning models, including the published baseline, when the building's input sensors are incomplete?}

\textbf{Sub-questions.}
\begin{enumerate}[label=RQ\arabic*,nosep]
\item How does accuracy change as 10--40\% of the sensor readings are removed?
\item Does the answer differ between random point loss (MCAR) and contiguous multi-sensor outages?
\item Does the physics term reduce physical inconsistency (balance residual, implausible responses), and are the learned thermal parameters credible?
\end{enumerate}

\textbf{Objectives.}
\begin{enumerate}[label=O\arabic*,nosep]
\item Build a reproducible benchmark with a fixed chronological split and identical, seeded missingness masks for every model.
\item Re-implement the baseline BNN and implement strong conventional baselines (profile, XGBoost, LSTM, GRU) and the proposed physics-informed GRU (PI-GRU) with architecture-matched ablations.
\item Evaluate accuracy, robustness, physical consistency and computational cost with five seeds, bootstrap confidence intervals and paired tests.
\item Beat the published baseline on its own protocol and state when the physics prior helps, is neutral, or hurts.
\end{enumerate}

\subsection{Contributions}
\begin{itemize}[nosep]
\item A sensor-loss benchmark for Building 59 with MCAR and contiguous outages at four rates, applied identically to every model, plus a meter-outage stress test.
\item A mask-aware PI-GRU that jointly predicts HVAC energy and indoor temperature and is regularised by a learnable RC heat balance with positive coefficients.
\item A faithful re-implementation of the baseline BNN, whose complete-data errors fall close to the published ones, and a direct comparison on the baseline's protocol.
\item Evidence, with confidence intervals and ablations, on where the accuracy gain comes from: the temporal, mask-aware architecture, the auxiliary temperature task or the physics term.
\item An open, one-command pipeline and a notebook that runs on Colab, Jupyter, macOS and Windows with CUDA, Apple-silicon MPS or CPU.
\end{itemize}

\subsection{Structure of the report}
Section~\ref{sec:lit} reviews related work and identifies the gap. Section~\ref{sec:method} describes the research design, data and evaluation protocol. Section~\ref{sec:design} specifies the models, and Section~\ref{sec:impl} the implementation. Section~\ref{sec:eval} presents the results and Section~\ref{sec:disc} discusses them. Section~\ref{sec:concl} concludes.
